\section{Related work}

The closest stationary-trajectory comparisons are \citet[Theorem 2.1, Corollary 2.2, and Theorem 3.1]{HuWager2023} and \citet[Theorem 1 (Uniform overlap frontier) and the Signed-depth family definition]{CausalSmithLatentOverlap2026}. Hu and Wager study stationary target-policy rewards under sequential ignorability, policy overlap, and mixing of the behavior and target dynamics. Their upper bounds accommodate multiple trajectories and heterogeneous units; their minimax lower bound concerns the stationary value of one unit observed through a single trajectory of length \leanref{sym:T}{\ensuremath{T}}, the number of observed periods. At fixed mixing scale \leanref{sym:t_0}{\ensuremath{t_0}} and logarithmic action likelihood-ratio bound \leanref{sym:\zeta}{\ensuremath{\zeta}}, their single-trajectory mean-squared-error rate is \(T^{-2/(2+t_0\zeta)}\). Partial-history importance weighting balances geometric forgetting against the variance cost of multiplying policy ratios. The CausalSmith source is a nonarchival public manuscript, version CSWP-2026-021 at the stable URL in the bibliography. Its cited conclusion enters the comparison as an external manuscript premise, with a trust status distinct from peer-reviewed evidence or an independently formalized dependency of this paper.
% [Hu and Wager, cited results](https://arxiv.org/pdf/2110.12343), [CausalSmith, cited manuscript](https://causalsmith.org/papers/stat_pomdp_latent_overlap_minimax_v1/paper.pdf).

Hidden-state cardinality makes the comparison precise. Both cited minimax experiments permit hidden alphabets whose size increases with trajectory length. The CausalSmith manuscript studies a \leanref{S-4}{cardinality-uniform bounded-stationary-overlap comparator}, with \leanref{sym:C}{\ensuremath{C}} bounding the ratio of target to behavior stationary probabilities on the joint state space. For fixed \(C>1\), \(t_0>0\), and \(\zeta>0\), its cited characterization has the same power rate. Its signed-depth construction uses \(2(Q+1)\) hidden states, where \leanref{sym:Q}{\ensuremath{Q}} is the terminal depth and increases logarithmically with trajectory length. Hu and Wager's lower construction likewise increases the hidden depth. Here, the \leanref{S-1}{single stationary behavior trajectory} comes from an experiment with binary observed states, binary hidden states, and binary actions. Under the law conditions in \cref{obj:def:binary-pomdp-class}, \cref{obj:thm:stable-grid-upper,obj:thm:fixed-binary-minimax} establish the matched \(T^{-1}\) mean-squared-error order for this fixed binary experiment. The different rates therefore concern different cardinality restrictions on the statistical experiment.

A second comparison concerns the information used to identify policy values. \citet[Section 4, Assumptions 2--4, and Theorems 2--4]{NairJiang2021} derive evaluation formulas from finite-horizon behavior-policy trajectories in a setting where action assignment may depend on the hidden state. Their spectral approach uses observable probability matrices whose rank equals the hidden-state cardinality. Extending the past enlarges the available proxy information; the future-based formula additionally uses distinct hidden-state reward probabilities \citep[Assumption 5 and Theorem 4]{NairJiang2021}. These conditions support recovery of the intervention distributions entering finite-horizon evaluation. In the present experiment, known action randomization conditional on the observed state, as specified in \cref{obj:ass:sequential-ignorability}, identifies the weighted intervention moments in \cref{obj:lem:observable-intervention-moments}. Uniform stability of the stationary scalar value then follows from \cref{obj:lem:pair-polynomial-modulus} over contracting four-state realizations, including realizations with reduced moment rank.
% [Nair and Jiang, identification results](https://arxiv.org/pdf/2109.10502).

The sampling and estimand distinctions also matter for \citet[Section 4 and Theorems 7--8]{ShiUeharaHuangJiang2022}. Their time-homogeneous analysis targets an infinite-horizon discounted value from a specified initial distribution and develops value and weight bridge functions for hidden-state-dependent behavior policies. The statistical analysis in Section 4 uses independent copies of observable transition tuples containing past and future observations. Theorem 7 establishes asymptotic normality under existence and uniqueness of the relevant bridge functions, boundedness, consistency, and a product condition involving weight-bridge error and the value-bridge conditional-moment residual. Theorem 8 gives a tabular identification formula for the single-decision specialization using a generalized inverse of an observable proxy matrix. These results connect informative proxies and bridge estimation to policy evaluation. The present analysis connects observed-state randomization and a fixed realization dimension to a uniform mean-squared-error bound for stationary mean reward from one dependent trajectory.
% [Shi, Uehara, Huang, and Jiang, cited results](https://proceedings.mlr.press/v162/shi22f/shi22f.pdf).

The finite-moment mechanism has a classical realization-theoretic interpretation. Finite-dimensional input-output representations and finite-rank sequence representations underlie the constructions of \citet[pp.\ 545--548]{Ho1966} and \citet[pp.\ 26--40]{Carlyle1971}. Here, the intervention moments admit \leanref{S-2}{four-state scalar linear realizations}, and \cref{obj:lem:pair-polynomial-modulus} controls the difference between stationary rewards by the first seven scalar moments. Contraction supplies uniform control of the stationary functional. This places the estimation target alongside, and distinguishes it from, richer learning objectives. \citet[Condition 1 and Theorem 6]{HsuKakadeZhang2012} estimate observable sequence distributions through a spectral representation, with sample requirements depending on singular values of emission and observable moment matrices. For latent-parameter estimation, \citet[Theorem 1 and Corollary 1]{AbrahamGassiatNaulet2023} quantify how learning two-state hidden Markov models depends on transition nondegeneracy and emission separation. The stability established here concerns the stationary reward functional across the contracting realization class; its uniform modulus is the property used by the estimator in \cref{obj:def:stable-grid-estimator}.
% [Carlyle and Paz, finite-rank realizations](https://www.sciencedirect.com/science/article/pii/S0022000071800053), [Hsu, Kakade, and Zhang, spectral learning](https://www.cs.columbia.edu/~djhsu/papers/hmm.pdf), [Abraham, Gassiat, and Naulet, parameter-learning bounds](https://arxiv.org/pdf/2106.12936).

Observable-value methods provide a closer functional comparison. \citet[Theorem 3]{Uehara2023} obtain a discounted-value error of order \(n^{-1/2}\), with instrument-strength and distribution-shift factors, under realization, Bellman completeness, latent-error identifiability, and coverage conditions. \citet[Theorem 7]{Zhang2024} give a finite-horizon bound of order \(H^2/\sqrt n\), up to class-complexity and outcome- and belief-coverage factors. Those results use independent finite-horizon samples and target discounted or finite-horizon values. Here the estimand is a stationary mean from one dependent trajectory, and contraction supplies uniform value stability across reduced-rank four-state realizations without a singular-value lower bound.

Fully observed stationary models provide useful parametric benchmarks with different estimands and nuisance requirements. \citet[Section 4 and Theorem 12]{KallusUehara2022} establish efficient evaluation of a normalized discounted policy value from stationary trajectory data. Their single-trajectory result uses sufficiently fast mixing, bounded action ratios, and a bounded ratio of target discounted state occupancy to the behavior stationary distribution. Cross-time fitting combines consistent estimates of the discounted action-value function and occupancy ratio, with the product of their \(L^2\) errors of order \(o(T^{-1/2})\). These conditions yield a parametric asymptotic scale for the discounted functional.
% [Kallus and Uehara, stationary-trajectory evaluation](https://arxiv.org/pdf/1909.05850).

For stationary mean reward, \citet[Section 3.2, Assumptions 7--8, and Theorem 3.9]{Mehrabi2024} analyze truncated doubly robust evaluation under weak distributional overlap. With bounded action ratios and their sampling, mixing, and boundedness conditions, a stationary occupancy ratio having a power-tail bound of order \(1+\eta\), where \(\eta>0\) is the tail parameter, gives stochastic mean-squared-error scales \(T^{-2\eta/(1+\eta)}\) for \(0<\eta<1\), \((\log T)/T\) for \(\eta=1\), and \(T^{-1}\) for \(\eta>1\). Their bound also includes the squared product of the \(L^2\) estimation errors for the differential action-value function and stationary occupancy ratio. This benchmark separates action overlap from stationary occupancy overlap and makes the nuisance contribution explicit. Against these comparisons, the contribution of \cref{obj:thm:stable-grid-upper,obj:thm:fixed-binary-minimax} is the particular matched characterization obtained by combining observable intervention moments, uniform scalar-value stability, and a lower bound within the same fixed binary class.
% [Mehrabi and Wager, weak-overlap evaluation](https://arxiv.org/pdf/2402.08201).
